\input{Iteration-V2/Chapter3/chapter-03-23}

\subsection{Revision 24 requirements and trust-store threat model}

Revision 24 adds eight requirements. \textbf{TS1} requires a fixed canonical schema whose anchor and revocation ordering cannot alter the signed payload. \textbf{TS2} requires one structurally valid Ed25519 offline root that is valid both at envelope issuance and assessment. \textbf{TS3} forbids an offline-root identifier from also identifying a reviewer anchor and requires different issuer and approver identities. \textbf{TS4} requires explicit, dated revocations within the envelope validity window. \textbf{TS5} requires sequence one to have no predecessor and later sequences to advance exactly one step while binding the accepted predecessor digest. \textbf{TS6} rejects rollback, same-sequence forks, gaps, missing history, invalid signatures, freeze through expiry, and future activation. \textbf{TS7} persists the highest accepted sequence and digest before exposing reviewer anchors and removes all effective anchors if persistence fails. \textbf{TS8} requires the interface to expose both the production provisioning state and local-history limitations.

The attacker model includes mutation of signed fields, substitution of an unknown or malformed root, replay of an older valid envelope, same-sequence replacement, omission of intermediate envelopes, predecessor substitution, and continued use after expiry. The design does not resist an attacker who can replace application code or roots, compromise the device, erase application storage, control the release channel, or subvert the offline ceremony. Those are residual distribution and platform threats, not cases silently counted as passing.
